% \usepackage{babel} 
% \usepackage{bidi} 
% \selectlanguage{farsi}




\chapter{مروری بر مطالعات انجام شده}
\section{مقدمه}

سیستم‌های حمل و نقل هوشمند \lr{ITS}\LTRfootnote{\lr{ITS} (\lr{Intelligent Transportation Systems})} به عنوان راهکاری برای توسعه روش‌های حمل و نقل کارآمد، ایمن و سازگار با محیط زیست شکل گرفته‌اند. این سیستم‌ها با ترکیب حسگرهای پیشرفته، شبکه‌های ارتباطی، تحلیل داده‌ها و الگوریتم‌های هوشمند طراحی شده‌اند. در مواجهه با چالش‌های پیچیده ناشی از شهرنشینی مدرن، این سیستم‌ها ضروری بوده و راه‌حل‌های نوآورانه‌ای برای کاهش ترافیک، بهبود امنیت و کاهش اثرات منفی زیست‌محیطی ارائه می‌دهند. تشخیص علائم راهنمایی و رانندگی، به عنوان یکی از حوزه‌های کلیدی در سیستم‌های حمل و نقل هوشمند، نقش بسیار مهمی در ارتقای ایمنی جاده‌ها، توسعه خودروهای خودران و بهبود سیستم‌های کنترل ترافیک ایفا می‌کند. این قابلیت می‌تواند به طور قابل‌توجهی بر ایمنی جاده‌ها و مدیریت ترافیک تأثیرگذار باشد. تشخیص دقیق و به موقع علائم، نقش اساسی در آگاه‌سازی رانندگان از محدودیت‌های سرعت و الزامات قانونی ایفا کرده و در نتیجه ایمنی جاده‌ها را به حداکثر می‌رساند. در رابطه با این موضوع رویکردهای گوناگونی برای مطالعه استفاده شده است [\ref{ref:Pagano}].

\section{روش‌های اولیه}

روش‌های اولیه عمدتاً بر تحلیل ویژگی‌های رنگ و شکل علائم متمرکز بودند. به‌عنوان مثال، پیکچولی و همکاران [\ref{ref:Piccioli}] (1996) با استفاده از رنگ‌های مشخص\LTRfootnote{Red, Green, Blue (\lr{RGB})}، نواحی کاندید را شناسایی کرده و از تطبیق الگو برای طبقه‌بندی استفاده می‌کردند. همچنین، زاده، آ. و کاسواند، ب. در 1998 [\ref{ref:Zadeh}] برای مقاوم‌سازی سیستم در برابر تغییرات مقیاس و زاویه، از تبدیل‌های هندسی\LTRfootnote{Fourier Transform Basis} بهره گرفتند. با این حال، این روش‌ها محدودیت‌هایی داشتند و نمی‌توانستند به‌خوبی با انسداد جزئی و شرایط نوری متغیر کنار بیایند.

برای بهبود عملکرد، روش‌های مبتنی بر شبکه عصبی\LTRfootnote{Neural Networks} معرفی شدند. کلمایر، و. و زواهلن [\ref{ref:Ellmeyer}] (1994) از اولین مقالاتی بود که شبکه‌های عصبی را برای تشخیص علائم به کار گرفت. این روش قابلیت یادگیری و تعمیم از داده‌های آموزشی را داشت اما در مواجهه با انسداد و نویز\LTRfootnote{Noise} عملکرد محدودی داشت. در ادامه، اسکالرا و همکاران [\ref{ref:Escalera}] (2003) روشی ترکیبی را معرفی کردند که از الگوریتم ژنتیک\LTRfootnote{Genetic Algorithm} برای شناسایی و شبکه‌های عصبی برای طبقه‌بندی استفاده می‌کرد. این روش با مقاومت در برابر تغییرات نور، انسداد جزئی و تغییر شکل علائم، توانست عملکرد بهتری در محیط‌های پیچیده ارائه دهد.
در سال 2005، بهلمان و همکاران [\ref{ref:Bahlmann}] یک سیستم جامع و سه‌مرحله‌ای برای شناسایی، ردیابی و طبقه‌بندی علائم ارائه کردند. در این روش، ابتدا با استفاده از الگوریتم تقویت تطبیقی\LTRfootnote{Adaptive Boosting (\lr{AdaBoost})}، ترکیبی از ویژگی‌های رنگی کانال‌های \lr{RGB}\LTRfootnote{Red, Green, Blue} و نرمال‌شده، ویژگی‌های هندسی (موقعیت، عرض و ارتفاع)، و دیگر پارامترها برای انتخاب قوی‌ترین ویژگی‌ها استفاده شد. سپس، با بهره‌گیری از یک مدل حرکتی ساده و ترکیب اطلاعات زمانی، نواحی شناسایی‌شده ردیابی شدند. در نهایت، این نواحی نرمال‌سازی شده و طبقه‌بندی با مدل‌های بیزین\LTRfootnote{Bayesian Models} انجام گرفت. این سیستم که روی 4000 تصویر آزمایش شد، توانست نرخ طبقه‌بندی 94٪ را بر روی داده‌های آزمون شامل 1700 تصویر به دست آورد. روش بهلمان با افزودن مرحله ردیابی حرکتی و استفاده از مدل‌های بیزین، پایداری و دقت بیشتری را ارائه داد. در حالی که روش اکاسکلرا [\ref{ref:Escalera}] به ترکیب الگوریتم ژنتیک و شبکه‌های عصبی تکیه داشت، استفاده از تقویت تطبیقی برای انتخاب ویژگی‌ها و مدل‌های بیزین برای طبقه‌بندی، قابلیت بهتری در مواجهه با مجموعه‌داده‌های بزرگ و متنوع فراهم کرد.
در مقاله گائو و همکاران [\ref{ref:Gao}] (2006)، از دو مدل \lr{CIECAM97} و \lr{FOSTS} برای شناسایی علائم راهنمایی و رانندگی استفاده شد. مدل \lr{CIECAM97} اطلاعات رنگی تصاویر را با تبدیل فضای \lr{RGB}\LTRfootnote{Red, Green, Blue} به \lr{HCL}\LTRfootnote{Hue, Chrome, Light} استخراج می‌کرد، در حالی که مدل \lr{FOSTS} با تقلید از سیستم بینایی انسان، ویژگی‌های شکلی علائم را از طریق تحلیل لبه‌ها و جهت‌گیری‌های تصویر شناسایی می‌کرد. این مدل با استفاده از ساختاری شامل دایره‌های هم‌مرکز و خطوط شعاعی، توانست ویژگی‌های شکلی را با مقاومت در برابر تغییرات مقیاس، چرخش، و تحریف استخراج کند. این سیستم با نرخ شناسایی 95٪ برای 98 تصویر علائم راهنمایی و رانندگی بریتانیایی، عملکرد مطلوبی حتی در شرایط نوری و زاویه‌ای متغیر ارائه داد.

در مقایسه، مالدونادو باسکان و همکاران [\ref{ref:Maldonado}] (2010) سیستمی مبتنی بر ماشین‌های بردار پشتیبان\LTRfootnote{Support Vector Machines (\lr{SVM})} توسعه دادند که دقت 95.5٪ را روی پایگاه داده‌ای شامل 36000 نمونه از 193 کلاس علائم اسپانیایی به دست آورد. این سیستم شامل سه مرحله اصلی بود: قطعه‌بندی علائم با ماسک‌های رنگی در فضای \lr{HSI}\LTRfootnote{Hue, Saturation, Intensity}، تشخیص شکل‌ها با تحلیل \lr{FFT}\LTRfootnote{Fast Fourier Transform} از لبه‌ها، و طبقه‌بندی علائم با ماشین‌های بردار پشتیبان. پیش‌پردازش‌هایی نظیر نرمال‌سازی خاکستری و افزایش کنتراست نیز دقت سیستم را 3-5٪ بهبود داد و گروه‌بندی بردارهای پشتیبان تعداد آن‌ها را تا 70٪ کاهش داد، که باعث افزایش کارایی سیستم شد.
در همین زمینه، هو و همکاران [\ref{ref:Hu}] (2010) از رویکردی مبتنی بر ویژگی‌های جزئی و مدل کیسه کلمات\LTRfootnote{Bag-of-Words (\lr{BoW})} استفاده کردند. در این روش، نقاط کلیدی تصویر با الگوریتم تبدیل ویژگی‌های مقیاس‌ناپذیر و چرخش‌ناپذیر\LTRfootnote{Scale-Invariant Feature Transform (\lr{SIFT})} شناسایی و توصیف شدند. ویژگی‌های \lr{SIFT} به دلیل مقاومت در برابر تغییرات مقیاس، چرخش، زاویه دید و نورپردازی انتخاب شدند. سپس این ویژگی‌ها با الگوریتم خوشه‌بندی\LTRfootnote{K-Means Clustering} و به فضای کیسه کلمات تبدیل شدند. در نهایت، طبقه‌بندی با ماشین بردار پشتیبان\LTRfootnote{Support Vector Machine (\lr{SVM})} و هسته تابع پایه شعاعی\LTRfootnote{Radial Basis Function (\lr{RBF})} انجام شد و پارامترهای $C$ و $\Gamma$ بهینه‌سازی گردیدند. این روش با دقت 93٪ روی 130 تصویر واقعی و زمان پردازش 0.098 ثانیه برای هر تصویر، عملکردی سریع و مقاوم در برابر تغییرات محیطی ارائه داد. در مقایسه، روش مالدونادو باسکان [\ref{ref:Maldonado}] به دلیل استفاده از داده‌های گسترده‌تر و پیش‌پردازش پیشرفته‌تر دقت بیشتری داشت، اما روش هو به دلیل زمان پردازش کمتر و انعطاف‌پذیری در شرایط متغیر برتری نشان داد. همچنین، روش گائو [\ref{ref:Gao}] با تمرکز بر مدل‌های بینایی انسان و مقاومت در برابر تغییرات، رویکرد متفاوتی ارائه کرد.

اما اکثر مدل‌هایی که نام برده شد در واقع برای مقایسه به طور کامل قابل اعتماد نیستند زیرا بر اساس داده‌های اختصاصی ارزیابی شده‌اند که در دسترس عموم نیستند. در سال 2011، استالکمپ مجموعه داده مرجع شناسایی علائم راهنمایی و رانندگی آلمان\LTRfootnote{German Traffic Sign Recognition Benchmark (\lr{GTSRB})} را معرفی کرد و به دنبال آن رقابت \lr{IJCNN} معرفی شد که در آن مدل‌های مختلفی بر روی این مجموعه داده ارزیابی شد [\ref{ref:Stallkamp}].
\section{شبکه های عصبی و یادگیری عمیق}
در یکی از نخستین نتایج موفقیت‌آمیز برای این دیتاستسرمنه، پ. و لکون، ی. (2011) [\ref{ref:Sermanet}]. ، از شبکه‌های عصبی کانولوشنی چندمقیاسی\LTRfootnote{Multi-Scale Convolutional Networks (\lr{CNNs})} با معماری پیشرفته استفاده شد که ویژگی‌های جزئی و کلی را از چندین لایه استخراج و برای طبقه‌بندی ترکیب می‌کرد. این معماری شامل چندین لایه کانولوشن، لایه‌های غیرخطی (مانند \lr{ReLU})، و لایه‌های \lr{Pooling} بود. هر لایه کانولوشن با استفاده از فیلترهایی با اندازه مشخص، ویژگی‌هایی مانند لبه‌ها، بافت‌ها و الگوهای پیچیده‌تر را استخراج می‌کرد. لایه‌های \lr{Pooling} برای کاهش ابعاد و افزایش مقاومت در برابر تغییرات جزئی مانند جابه‌جایی و چرخش طراحی شده بودند. علاوه بر این، خروجی‌های چندین لایه استخراج ویژگی به صورت هم‌زمان به طبقه‌بند نهایی وارد می‌شدند که ترکیبی از اطلاعات جزئی و کلی را فراهم می‌کرد.

برای افزایش تنوع داده‌ها، از تکنیک‌های تقویت داده\LTRfootnote{Data Augmentation} مانند چرخش، جابه‌جایی، تغییر مقیاس و افزودن نویز استفاده شد. این تغییرات به‌طور تصادفی در هر دوره آموزشی اعمال می‌شدند و باعث افزایش پنج‌برابری نمونه‌های آموزشی و بهبود مقاومت مدل در برابر تغییرات محیطی شدند. این روش توانست دقتی نزدیک به 99٪ را به‌دست آورد و یکی از موفق‌ترین کاربردهای اولیه شبکه‌های عصبی کانولوشنی در این حوزه محسوب می‌شد.

در ادامه این تلاش‌ها، یکی از برجسته‌ترین این روش‌ها، شبکه‌های عصبی چندستونه\LTRfootnote{Multi-Column Deep Neural Networks (\lr{MCDNN})} ارائه‌شده توسط تیم \lr{IDSIA} بود [\ref{ref:Ciresan}]. در این روش، چندین شبکه عصبی کانولوشنی به‌صورت موازی ترکیب شدند تا عملکرد طبقه‌بندی بهبود یابد. این سیستم با استفاده از داده‌های خام و تقویت داده‌ها\LTRfootnote{Data Augmentation}، به‌طور هم‌زمان ویژگی‌های سطح پایین و سطح بالا را استخراج می‌کرد. در این رویکرد، اعمال تغییرات تصادفی بر روی داده‌ها در هر دوره آموزشی منجر به افزایش تعمیم‌پذیری مدل شد. این روش با دقت 99.46٪، عملکردی فراتر از توانایی انسانی ارائه داد. علاوه بر این، جنگل تصادفی\LTRfootnote{Random Forest (\lr{RF})} [\ref{ref:RF}] و تحلیل تفکیکی خطی\LTRfootnote{Linear Discriminant Analysis (\lr{LDA})} [\ref{ref:LDA}] نیز برای مقایسه با روش‌های پیچیده‌تر به کار گرفته شدند. جنگل تصادفی با استفاده از ترکیب چندین درخت تصمیم که بر زیرمجموعه‌های مختلف داده‌ها آموزش داده شده بودند، عملکردی پایدار در برابر عدم توازن داده‌ها ارائه کرد. تحلیل تفکیکی خطی، علیرغم سادگی و هزینه محاسباتی پایین، با استفاده از ویژگی‌های هیستوگرام گرادیان‌های جهت‌دار\LTRfootnote{Histogram of Oriented Gradients (\lr{HOG})} [\ref{ref:HOG}] توانست دقتی بالاتر از 95٪ کسب کند. این نشان می‌داد که حتی مدل‌های کم‌هزینه نیز می‌توانند در شرایط خاص رقابتی باشند.

شبکه عصبی کانولوشنی سریع\LTRfootnote{Faster Region-Based Convolutional Neural Network (\lr{Faster R-CNN})} [\ref{ref:Ren}]، به‌عنوان یکی از پیشرفته‌ترین مدل‌های تشخیص اشیاء، در سال 2015 توسط رِن و همکاران ارائه شد. این مدل با ترکیب شبکه پیشنهاد منطقه\LTRfootnote{Region Proposal Network (\lr{RPN})} و شبکه عصبی کانولوشنی در یک چارچوب یکپارچه، امکان تشخیص سریع و دقیق را فراهم کرد. در این مدل، شبکه پیشنهاد منطقه با اسکن تصویر و تولید نواحی پیشنهادی با احتمال بالاتر حضور اشیاء، نقش کلیدی ایفا می‌کند. سپس این نواحی به شبکه عصبی کانولوشنی اصلی ارسال می‌شوند تا ویژگی‌ها استخراج و اشیاء طبقه‌بندی شوند. شبکه عصبی کانولوشنی سریع توانست مشکلات روش‌های قبلی مانند وابستگی به الگوریتم‌های خارجی برای تولید نواحی پیشنهادی را برطرف کرده و دقت و سرعت بیشتری ارائه دهد.

به الگوریتم‌های خارجی برای تولید نواحی پیشنهادی را برطرف کرده و دقت و سرعت بیشتری ارائه دهد.  
در این پژوهش نیز [\ref{ref:Zhu}] از \lr{Faster R-CNN} برای تشخیص علائم راهنمایی استفاده شده است. در این روش، \lr{ResNet50} به‌عنوان شبکه استخراج ویژگی انتخاب شده، زیرا لایه‌های عمیق آن قابلیت استخراج ویژگی‌های پیچیده و دقیق را دارند. همچنین، جعبه‌های لنگر\LTRfootnote{Anchor Boxes} برای نواحی پیشنهادی بهینه‌سازی شده‌اند تا با اندازه و شکل کوچک علائم راهنمایی سازگاری بیشتری داشته باشند. فرآیند شامل شناسایی نواحی مرتبط توسط شبکه پیشنهاد منطقه، استخراج ویژگی‌ها با شبکه عصبی کانولوشنی و طبقه‌بندی نهایی علائم به کلاس‌های مختلف است.  

در الگوریتم بهبود‌یافته تشخیص علائم راهنمایی بر اساس \lr{Faster R-CNN}، استفاده از کانولوشن‌های تغییرپذیر\LTRfootnote{Deformable Convolutions} و هرم ویژگی‌ها\LTRfootnote{Feature Pyramid Network (\lr{FPN})} [\ref{ref:Zhu}] از نکات برجسته است. کانولوشن‌های تغییرپذیر به دلیل انعطاف‌پذیری بالا در مواجهه با تغییرات شکل و مقیاس علائم، فرآیند استخراج ویژگی‌ها را بهبود می‌دهند. هرم ویژگی نیز با ترکیب اطلاعات از لایه‌های مختلف شبکه، امکان استخراج ویژگی‌های دقیق‌تر و سطح بالا از علائم کوچک و پیچیده را فراهم می‌کند.  

در پژوهش دیگری [\ref{ref:Li}] از نسخه بهینه‌شده \lr{ResNet50-D} برای استخراج ویژگی‌های قوی‌تر استفاده شده است. همچنین، ماژول توجه به هرم ویژگی‌ها اضافه شده است تا تمرکز بیشتری روی نواحی مهم مانند علائم جاده‌ای ایجاد شود. این ماژول توجه، نواحی پیشنهادی تولید‌شده را تقویت کرده و به طبقه‌بندی‌کننده ارسال می‌کند که دقت سیستم را در محیط‌های شلوغ و پیچیده جاده‌ای افزایش می‌دهد.  

نتایج این مطالعات نشان داده‌اند که با استفاده از بهبودهایی مانند کانولوشن‌های تغییرپذیر، ماژول‌های توجه و بهینه‌سازی جعبه‌های لنگر می‌توان \lr{Faster R-CNN} را برای کاربردهای خاص مانند تشخیص علائم راهنمایی بهینه کرد. این بهبودها باعث افزایش مقاومت مدل در برابر تغییرات محیطی و ارائه نتایج سریع‌تر و دقیق‌تر می‌شوند.  
مدل \lr{YOLO}\LTRfootnote{You Only Look Once} [\ref{ref:Redmon}], که در سال 2016 توسط جوزف ردمن و همکاران معرفی شد، تحولی اساسی در حوزه شناسایی اشیاء ایجاد کرد و به‌ویژه در تشخیص علائم راهنمایی و رانندگی کاربردهای چشمگیری داشت. برخلاف روش‌های سنتی نظیر \lr{R-CNN} و \lr{Faster R-CNN} که فرآیند شناسایی را در دو مرحله جداگانه انجام می‌دادند، \lr{YOLO} توانست این مراحل را در یک گام یکپارچه کند. این تغییر، سرعت و کارایی مدل را به شکل چشمگیری افزایش داد. نسخه اولیه این مدل\LTRfootnote{\lr{YOLOv1}} [\ref{ref:Redmon}] با تأکید بر سرعت، امکان شناسایی بلادرنگ اشیاء را فراهم کرد، اما محدودیت‌هایی در شناسایی اشیاء کوچک داشت. نسخه‌های بعدی مانند \lr{YOLOv2} و \lr{YOLOv3} با بهبود معماری شبکه و تقویت قابلیت تشخیص اشیاء در مقیاس‌های مختلف، این نقاط ضعف را کاهش دادند. \lr{YOLOv3} با ارائه تعادلی مناسب بین دقت و سرعت، به‌عنوان پایه‌ای برای توسعه نسخه‌های پیشرفته‌تر از جمله \lr{YOLOv4}، \lr{YOLOv5}، و \lr{YOLOv7} شناخته شد.

کاربرد \lr{YOLO} در تشخیص علائم راهنمایی و رانندگی با توسعه نسخه‌های جدید گسترش یافت. در یکی از تحقیقات جدید، مدل \lr{YOLOv7} در مقاله‌ای به نام \lr{YOLO-CCA} توسط جیانگ و همکاران (2024) [\ref{ref:Jiang}] برای شناسایی علائم راهنمایی مورد استفاده قرار گرفت. این مدل با افزودن ماژول‌هایی برای تقویت ویژگی‌های جزئی و جمع‌آوری اطلاعات کلی، توانست دقت تشخیص علائم پیچیده را بهبود بخشد و عملکرد بهتری در شرایط متغیر محیطی ارائه دهد. این تحقیق با استفاده از مجموعه داده‌های واقعی شامل تصاویری از علائم در شرایط مختلف، مانند تغییرات نور و زاویه دید، انجام شد.

در مقاله دیگری با عنوان \lr{YOLO-TS} که توسط چن و همکاران در سال 2024 معرفی شد [\ref{ref:Chen}], بهینه‌سازی میدان‌های دریافت و حذف انکرها\LTRfootnote{Anchor-Free Detection} به‌منظور افزایش دقت و سرعت تشخیص علائم انجام شد. در این پژوهش از مجموعه داده‌های ترکیبی مانند \lr{GTSRB} و سایر مجموعه‌های مشابه استفاده شده است. این تغییرات با کاهش وابستگی به انکرهای دستی، انعطاف‌پذیری بیشتری برای شناسایی علائم در محیط‌های پویا ایجاد کردند.

تحقیقات دیگری نیز بر بهبود معماری \lr{YOLO} متمرکز بوده‌اند. در مقاله \lr{Improved YOLOv5} که توسط وانگ و همکاران در سال 2021 منتشر شد [\ref{ref:Wang}], تغییراتی نظیر افزودن هرم ویژگی توجهی\LTRfootnote{Attention Feature Pyramid} و ماژول تقویت ویژگی‌ها به \lr{YOLOv5} اعمال گردید. این تغییرات توانایی مدل را در شناسایی علائم کوچک و پیچیده بهبود بخشید و دقت آن را در شرایط محیطی متغیر افزایش داد. این پژوهش از مجموعه داده علائم راهنمایی و رانندگی \lr{LISA}\LTRfootnote{Laboratory for Intelligent and Safe Automobiles (\lr{LISA})} [\ref{ref:LISA}], که شامل تصاویر علائم راهنمایی آمریکایی با شرایط متنوع محیطی است، بهره برد.

مقاله دیگری به نام \lr{Sign-YOLO} نیز، که از \lr{YOLOv7} استفاده کرده است [\ref{ref:SignYOLO}], با بهره‌گیری از مکانیزم توجه توانست دقت شناسایی علائم پیچیده را به شکل چشمگیری افزایش دهد. مدل \lr{TRD-YOLO} که توسط \lr{MDPI Sensors} در سال 2023 معرفی شد، به‌طور ویژه برای شناسایی علائم کوچک طراحی شده است. این مدل با استفاده از ماژول استخراج ویژگی جهانی و سر تشخیص سبک، عملکرد فوق‌العاده‌ای در این زمینه ارائه داد. مجموعه داده \lr{STS} (Small Traffic Sign Dataset)، که به‌طور خاص برای علائم کوچک طراحی شده بود، به‌عنوان پایه داده این تحقیق مورد استفاده قرار گرفت.

نسخه‌های مختلف \lr{YOLO} مانند \lr{YOLOv3} و نسخه‌های جدیدتر نظیر \lr{YOLOv8} نیز تأثیر شایانی در پیشرفت تشخیص علائم راهنمایی و رانندگی داشته‌اند. \lr{YOLOv3} با پایداری بالا، توانایی شناسایی علائم رایج در کشورهای اروپایی را با استفاده از مجموعه داده \lr{GTSRB} فراهم کرد، در حالی که \lr{YOLOv8} با بهبود قابل‌توجه در دقت و سرعت، برای شناسایی علائم پیچیده‌تر مورد استفاده قرار گرفت. این نسخه‌ها با ارائه قابلیت‌هایی مانند حذف انکرها، تقویت ویژگی‌ها، و استفاده از مکانیزم‌های توجه، به ابزارهایی بسیار کارآمد برای شناسایی بلادرنگ علائم راهنمایی و رانندگی تبدیل شده‌اند و نقش مهمی در توسعه سیستم‌های هوشمند رانندگی و بهبود ایمنی جاده‌ای ایفا کرده‌اند.

\section{مدل های ترنسفورمر}
مدل‌های ترنسفورمر، که در سال 2017 معرفی شدند [\ref{ref:Vaswani}]، با استفاده از مکانیسم خودتوجهی\LTRfootnote{Self-Attention Mechanism} [\ref{ref:Vaswani}] و توانایی پردازش هم‌زمان داده‌ها، انقلابی در یادگیری عمیق ایجاد کردند. این مدل‌ها ابتدا در حوزه پردازش زبان طبیعی مورد استفاده قرار گرفتند، اما به‌سرعت توانستند به حوزه‌های بینایی کامپیوتر هم گسترش یابند. معرفی ترنسفورمرهای بینایی\LTRfootnote{Vision Transformers (\lr{ViT})} [\ref{ref:Dosovitskiy}] در سال 2020 نشان داد که این رویکرد می‌تواند در وظایف پیچیده بینایی نیز عملکرد بی‌نظیری داشته باشد. در ترنسفورمرهای بینایی، تصاویر به قطعات کوچکی تقسیم می‌شوند که هر قطعه به‌عنوان یک توکن جداگانه پردازش می‌شود. این توکن‌ها سپس با مکانیسم خودتوجهی بررسی می‌شوند تا ارتباطات میان بخش‌های مختلف تصویر تحلیل شود. این ساختار به مدل اجازه می‌دهد تا به‌طور هم‌زمان جزئیات دقیق و ویژگی‌های کلی تصویر را تحلیل کند و در وظایف مختلف بینایی، بهبود چشمگیری در دقت و کارایی ایجاد نماید.

یکی از مهم‌ترین پیشرفت‌ها در استفاده از ترنسفورمرها، توسعه مدل‌های شناسایی اشیاء مانند \lr{DETR}\LTRfootnote{Deformable Transformer (\lr{DETR})} [\ref{ref:Carion}] بود که به‌عنوان جایگزینی برای روش‌های سنتی مبتنی بر \lr{CNN} معرفی شد. مدل \lr{DETR} با حذف فرآیندهای چندمرحله‌ای مانند تولید نواحی پیشنهادی و استفاده مستقیم از خودتوجهی، توانست دقت بالایی را در شناسایی اشیاء با مقیاس‌های مختلف ارائه دهد. این ویژگی، به‌ویژه در زمینه شناسایی علائم راهنمایی و رانندگی، کاربرد قابل‌توجهی داشته است. مدل‌های ترنسفورمرهای بینایی و \lr{DETR} با توانایی خود در تحلیل روابط پیچیده میان عناصر تصویر، عملکرد بهتری در شناسایی علائم کوچک، علائم با شرایط نوری متفاوت و علائم پیچیده از خود نشان داده‌اند.

در پژوهش‌های اخیر، استفاده از ترنسفورمرهای بینایی و مدل‌های مبتنی بر ترنسفورمر در شناسایی علائم راهنمایی و رانندگی به نتایج درخشانی منجر شده است. به‌عنوان مثال، در [\ref{ref:Mingwin}] با استفاده از مجموعه داده‌های آلمان و بلژیک\LTRfootnote{Belgium Traffic Sign Dataset (\lr{BelgiumTS})} [\ref{ref:BelgiumTS}]، مدل توانست عملکردی بسیار دقیق در شناسایی علائم به نمایش بگذارد. همچنین، مدل ترنسفورمر بینایی محلی\LTRfootnote{Local Vision Transformer (\lr{LVT})} [\ref{ref:Farzinpour}] که ترکیبی از شبکه‌های عصبی کانولوشنی و ترنسفورمرها است و توسط فرزین پور و همکاران معرفی شده، با استفاده از ماژول محلی‌سازی توانسته است دقت 99.8٪ را در مجموعه داده راهنمایی و رانندگی ایران [\ref{ref:IranDataset}] و 99.66٪ را در مجموعه داده آلمان به دست آورد. این مدل با کاهش پیچیدگی محاسباتی و تمرکز بیشتر بر نواحی مرتبط با علائم، کارایی قابل‌توجهی از خود نشان داده است.

از سوی دیگر، مقاله [\ref{ref:Xia}] با استفاده از نسخه‌ای بهبودیافته از \lr{DETR}، شامل ماژول‌های پولینگ هرمی فضایی گسترش‌یافته\LTRfootnote{Dilated Spatial Pyramid Pooling (\lr{DSPP})} [\ref{ref:DSPP}] و ماژول تجمیع باقی‌مانده ویژگی‌ها\LTRfootnote{Feature Residual Aggregation Module (\lr{FRAM})} [\ref{ref:FRAM}] توانسته است عملکردی استثنایی در شناسایی علائم با مقیاس‌های مختلف ارائه دهد. این تحقیق با استفاده از مجموعه داده‌های آلمان و چین، دقت بالایی را در شناسایی علائم کوچک و بزرگ گزارش کرده است. همچنین در مقاله دیگری [\ref{ref:RTDETR}], نسخه‌ای سبک‌تر از \lr{RT-DETR} معرفی شده که با کاهش پیچیدگی محاسباتی، توانسته است شناسایی بلادرنگ علائم راهنمایی را بهبود بخشد. این مدل بر روی مجموعه داده آلمان آزمایش شده و عملکرد قابل‌توجهی را از نظر سرعت و دقت ارائه داده است.

یکی از ترکیب‌های نوآورانه، استفاده از مدل‌های ترنسفورمر بینایی کارآمد\LTRfootnote{Efficient Vision Transformers} [\ref{ref:EfficientViT}] همراه با \lr{YOLOv5} در پژوهشی بود که توانست از سرعت بالای \lr{YOLOv5} و دقت بالای ترنسفورمرها به‌طور هم‌زمان بهره‌مند شود. این ترکیب بر روی مجموعه داده آلمانی آزمایش شده و نتایج درخشانی در شناسایی علائم راهنمایی و رانندگی نشان داده است. این پیشرفت‌ها، به‌ویژه با استفاده از مجموعه داده‌های استاندارد نظیر آلمان، بلژیک، ایران، و چین، توانسته‌اند عملکرد مدل‌ها را در شرایط متنوع محیطی شامل تغییرات نوری، زوایای مختلف و مقیاس‌های متفاوت، بهبود دهند. مدل‌های ترنسفورمر با دقت بالا و انعطاف‌پذیری خود، ابزار کلیدی برای توسعه سیستم‌های هوشمند رانندگی و ارتقای ایمنی جاده‌ها هستند و به‌عنوان یکی از پایه‌های اصلی فناوری‌های نوین بینایی کامپیوتر شناخته می‌شوند.





\section{منابع}


\begin{LTR}






\begin{latin}
\begin{enumerate}





    \item \label{ref:Pagano}  , P. (2021). "Intelligent Transportation Systems: From Good Practices to Standards." CRC Press, Boca Raton, FL, USA.
    \item \label{ref:Piccioli} Piccioli, A., et al. (1996). "Identification of candidate regions using specific colors and classification using pattern matching."
    \item \label{ref:Zadeh} Zadeh, A., & Kasvand, B. (1998). "Geometric transformations for enhancing system robustness against scale and rotation changes."
    \item \label{ref:Ellmeyer} Ellmeyer, W., & Zwahlen, (1994). "The first papers that applied neural networks for sign recognition."
    \item \label{ref:Escalera} Escalera, S., & Others. (2003). "A hybrid method combining genetic algorithms for detection and neural networks for classification."
    \item \label{ref:Bahlmann} Bahlmann, C., Zhu, Y., Ramesh, V., Pellkofer, M., and Koehler, T. (2005). "A system for traffic sign detection, tracking, and recognition using color, shape, and motion information." In Proceedings of the IEEE Intelligent Vehicles Symposium, pages 255–260. IEEE Press.
    \item \label{ref:Gao} Gao, X. W., Podladchikova, L., Shaposhnikov, D., Hong, K., and Shevtsova, N. (2006). "Recognition of traffic signs based on their colour and shape features extracted using human vision models." Journal of Visual Communication and Image Representation, 17(4):675–685.
    \item \label{ref:Maldonado} Maldonado Bascon, S., et al. (2010). "Optimization on pictogram identification for road-sign recognition task using SVMs." Computer Vision and Image Understanding.
    \item \label{ref:Hu} Hu, et al. (2010). "Local feature-based approach and Bag-of-Words model for sign recognition."
    \item \label{ref:Stallkamp} Stallkamp, J., Schlipsing, M., Salmen, J., & Igel, C. (2011). "Man vs. Computer: Benchmarking Machine Learning Algorithms for Traffic Sign Recognition." 2011 IEEE International Joint Conference on Neural Networks (IJCNN), 3009–3016.
    \item \label{ref:Sermanet} Sermanet, P., & LeCun, Y. (2011). "Traffic Sign Recognition with Multi-Scale Convolutional Networks." In Proceedings of the International Joint Conference on Neural Networks (IJCNN), pp. 2809–2813.
    \item \label{ref:Ciresan} Ciresan, D. C., Meier, U., Masci, J., Gambardella, L. M., & Schmidhuber, J. (2012). "Multi-column deep neural networks for traffic sign classification." Neural Networks, 32, 333–338.
    \item \label{ref:RF} Breiman, L. (2001). "Random Forests." Machine Learning, 45(1), 5–32.
   \item \label{ref:LDA} Fisher, R. A. (1936). "The use of multiple measurements in taxonomic problems." Annals of Eugenics, 7(2), 179–188.

    \item \label{ref:HOG} Dalal, N., & Triggs, B. (2005). "Histograms of Oriented Gradients for Human Detection." In Proceedings of the IEEE Computer Society Conference on Computer Vision and Pattern Recognition (CVPR), 886–893.
    \item \label{ref:Ren} Ren, S., He, K., Girshick, R., & Sun, J. (2015). "Faster R-CNN: Towards Real-Time Object Detection with Region Proposal Networks." Advances in Neural Information Processing Systems (NeurIPS), 28, 91–99.
    
    \item \label{ref:Zhu} Zhu, H., & Li, D. (2020). "Traffic sign detection using Faster R-CNN with FPN and deformable convolutions."
    \item \label{ref:Li} Li, X., Xie, Z., Deng, X., Wu, Y., & Pi, Y. (2021). "Traffic sign detection based on improved faster R-CNN for autonomous driving." The Journal of Supercomputing, 78(2), 1458–1480.
        \item \label{ref:Redmon} Redmon, J., Divvala, S., Girshick, R., & Farhadi, A. (2016). "You Only Look Once: Unified, Real-Time Object Detection." Proceedings of the IEEE Conference on Computer Vision and Pattern Recognition (CVPR), 779–788.
    \item \label{ref:Jiang} Jiang, L., Zhan, P., Bai, T., & Yu, H. (2024). "YOLO-CCA: A Context-Based Approach for Traffic Sign Detection."
    \item \label{ref:Chen} Chen, J., Huang, H., Zhang, R., Lyu, N., Guo, Y., Dai, H.-N., & Yan, H. (2024). "YOLO-TS: Real-Time Traffic Sign Detection with Enhanced Accuracy Using Optimized Receptive Fields and Anchor-Free Fusion."
    \item \label{ref:Wang} Wang, J., Chen, Y., Gao, M., & Dong, Z. (2021). "Improved YOLOv5 network for real-time multi-scale traffic sign detection."
    \item \label{ref:LISA} LISA: Laboratory for Intelligent and Safe Automobiles. "LISA Traffic Sign Dataset."
    \item \label{ref:SignYOLO} Chu, J., Zhang, C., Yan, M., Zhang, H., & Ge, T. (2023). "TRD-YOLO: A Real-Time, High-Performance Small Traffic Sign Detection Algorithm." Sensors, 23(8), 3871.
    \item \label{ref:Vaswani} Vaswani, A., et al. (2017). "Attention Is All You Need." Advances in Neural Information Processing Systems.
    \item \label{ref:Dosovitskiy} Dosovitskiy, A., et al. (2020). "An Image is Worth 16x16 Words: Transformers for Image Recognition at Scale."
    \item \label{ref:Carion} Carion, N., et al. (2020). "End-to-End Object Detection with Transformers." In European Conference on Computer Vision.
    \item \label{ref:Mingwin} Mingwin, S., et al. (2024). "Revolutionizing Traffic Sign Recognition: Unveiling the Potential of Vision Transformers."
    \item \label{ref:BelgiumTS} Belgium Traffic Sign Dataset (BelgiumTS).
    \item \label{ref:Farzinpour} Farzinpour, M., Rahmani, H., & Ebrahimi, M. (2023). "Local Vision Transformer: Efficient Integration of CNN and Transformer for Traffic Sign Recognition." IEEE Transactions on Intelligent Transportation.
    \item \label{ref:IranDataset} Persian Traffic Sign Dataset.
    \item \label{ref:Xia} Xia, J., et al. (2023). "DSRA-DETR: An Improved DETR for Multiscale Traffic Sign Detection." Sustainability, 15(14), 10862.
    \item \label{ref:DSPP} Dilated Spatial Pyramid Pooling (DSPP).
    \item \label{ref:FRAM} Feature Residual Aggregation Module (FRAM).
    \item \label{ref:RTDETR} RT-DETR: Traffic Sign Detection Based on Lightweight Real-Time DETR.
    \item \label{ref:EfficientViT} Efficient Vision Transformers for Traffic Sign Recognition.
\end{enumerate}
\end{latin}

\end{LTR}





